% titel: Blatt 6, A5: Welche Verbindungen sind nicht bildbar?

\subsection*{Frage}
Welche der Verbindungen $\mathrm{NCl_3,\ NCl_5,\ PCl_3,\ PCl_5,\ CH_5,\ HClO_4,\ SeF_4}$
können nicht gebildet werden?

\subsection*{Lösung}
\textbf{Nicht bildbar:}
\begin{itemize}
  \item \textbf{$\mathrm{NCl_5}$:} N steht in Periode 2, hat keine $d$-Orbitale
        und kann sein Oktett nicht erweitern ($>8$ Valenzelektronen unmöglich).
  \item \textbf{$\mathrm{CH_5}$:} C kann höchstens 4 Bindungen bilden; in
        $\mathrm{CH_5}$ hätte C 5 Partner (beziehungsweise überzoges
        Oktett) --- nicht realisierbar.
\end{itemize}
\textbf{Bildbar:} $\mathrm{NCl_3}$, $\mathrm{PCl_3}$ (Oktett erfüllt),
$\mathrm{PCl_5}$ (P in Periode 3 kann Oktett erweitern),
$\mathrm{HClO_4}$ (Cl erweitert, formal hohe Valenz), $\mathrm{SeF_4}$
(Se in Periode 4 mit Oktetterweiterung).

\begin{ergebnis}
Nicht bildbar: $\mathrm{NCl_5}$ und $\mathrm{CH_5}$. Der Grund ist jeweils die
Oktettgrenze wirksamer 2.-Periode-Zentralatome (N, C).
\end{ergebnis}

\begin{kontrolle}
N vs. P: $\mathrm{NCl_3}<=\mathrm{PCl_3}$ (beide Oktett), aber
$\mathrm{NCl_5}$ scheitert, während $\mathrm{PCl_5}$ existiert --- exakt der
Periodeneffekt. $\mathrm{CH_5}$: C kann nie fünfbindig sein; auch die
H-Ungleichung der Elektronenbilanz passt nicht (C hätte $>8\,e^-$).
\end{kontrolle}

\begin{merksatz}
Ist das Zentralatom in Periode 2, gibt es keine Oktetterweiterung; C maximal 4,
$\mathrm{Cl^-}$, N maximal 3-4 Bindungen. P, S, Se, I dürfen erweitern.
\end{merksatz}
